% One-page CV template "onepage-blue" (v2 - bullets)
% Plain article class, Helvetica-style sans (TeX Gyre Heros), moderncv blue
% accents on the name, section headings, rules and employer lines. Single
% column, so the ATS text layer reads top to bottom in visual order.
% Compile: lualatex (fontspec). Hard limit: exactly 1 page.
\documentclass[10pt,letterpaper]{article}

\usepackage[top=0.6in,bottom=0.55in,left=0.7in,right=0.7in]{geometry}
\usepackage{fontspec}
\setmainfont{TeX Gyre Heros}
\setsansfont{TeX Gyre Heros}
\usepackage{xcolor}
% moderncv "blue" color1 (rgb 0.22,0.45,0.70)
\definecolor{cvblue}{rgb}{0.22,0.45,0.70}
\definecolor{cvgray}{rgb}{0.35,0.35,0.35}
\usepackage{titlesec}
\usepackage{enumitem}
\usepackage{needspace}
\usepackage[hidelinks]{hyperref}
\hypersetup{
    pdftitle={[YOUR_NAME] - CV},
    pdfauthor={[YOUR_NAME]},
}

\pagestyle{empty}
\setlength{\parindent}{0pt}
\setlength{\parskip}{0pt}
\linespread{1.08}
\emergencystretch=1em
\frenchspacing

% Section heading: bold uppercase in blue, full-width blue rule beneath.
\titleformat{\section}
    {\large\bfseries\color{cvblue}}
    {}{0pt}{\MakeUppercase}
    [\vspace{-0.4ex}{\color{cvblue}\titlerule[0.6pt]}]
\titlespacing*{\section}{0pt}{9pt}{5pt}

% \cventry{Role / degree}{Employer -- location}{Start - End}
% Line 1: bold role, date right-aligned. Line 2: employer in blue.
% Keep #1 short enough for one line; put team/location in #2, never in #1.
\newcommand{\cventry}[3]{%
    \needspace{5\baselineskip}%
    \par\noindent\textbf{#1}\hfill\mbox{#3}\par\nopagebreak
    \noindent{\color{cvblue}#2}\par\nopagebreak\vspace{1pt}}
% \begin{cvbullets} \item ... \end{cvbullets} - 3 to 4 bullets per role,
% each one printed line (two at most), verb first, a number where possible.
\newenvironment{cvbullets}
    {\begin{itemize}[leftmargin=1.2em,topsep=1pt,itemsep=1pt,parsep=0pt,label={\color{cvblue}\textbullet}]}
    {\end{itemize}\vspace{4pt}}
% \cvline{Bold title}{inline description}{date} - one-line items
\newcommand{\cvline}[3]{%
    \noindent\parbox[t]{0.89\textwidth}{\raggedright\textbf{#1} -- #2}%
    \hfill\parbox[t]{0.1\textwidth}{\raggedleft #3}\par\vspace{1pt}}
% \cvskill{Label}{comma-separated skills} - one line of the Skills section
\newcommand{\cvskill}[2]{\noindent\textbf{#1:} #2\par\vspace{1pt}}
% \cvcert{Certification}{Date}
\newcommand{\cvcert}[2]{\noindent\textbf{#1}\hfill #2\par}

\begin{document}

% ---------- Header ----------
{\fontsize{22}{26}\selectfont\bfseries\color{cvblue}\MakeUppercase{[YOUR_NAME]}\par}
\vspace{2pt}
{\large [TARGET_TITLE - the posting's role title or closest truthful equivalent]\par}
\vspace{4pt}
\textbf{Email:} \href{mailto:[YOUR_EMAIL]}{[YOUR_EMAIL]} \quad | \quad \textbf{Phone:} [YOUR_PHONE]\par
\textbf{GitHub:} \href{https://[YOUR_GITHUB]}{[YOUR_GITHUB]} \quad | \quad \textbf{LinkedIn:} \href{https://[YOUR_LINKEDIN]}{[YOUR_LINKEDIN]}\par

% ---------- Summary ----------
\section{Professional Summary}
[PROFILE_STATEMENT - 3 lines max: who you are, the 2-3 frameworks/tools the posting asks for, one headline result. Do not repeat the bullets.]\par

% ---------- Skills ----------
\section{Skills}
\cvskill{[LABEL]}{[items]}
\cvskill{[LABEL]}{[items]}
\cvskill{[LABEL]}{[items]}

% ---------- Experience ----------
\section{Experience}
\cventry{[ROLE]}{[COMPANY] -- [TEAM] -- [LOCATION]}{[START] - [END]}
\begin{cvbullets}
    \item [Verb + what + result, posting's own terms]
    \item [...]
    \item [...]
\end{cvbullets}

\cventry{[ROLE]}{[COMPANY] -- [TEAM] -- [LOCATION]}{[START] - [END]}
\begin{cvbullets}
    \item [...]
    \item [...]
    \item [...]
    \item [...]
\end{cvbullets}

% ---------- Education ----------
\section{Education}
\cventry{[DEGREE]}{[UNIVERSITY], [CITY]}{[START] - [END]}
[One line: thesis, honours - no GPA]\par

% ---------- Projects (optional - only if relevant to the posting) ----------
\section{Projects}
\cvline{[SHORT PROJECT]}{[one-line description]}{[YEAR]}

% ---------- Certifications ----------
\section{Courses and Certifications}
\cvcert{[CERTIFICATION]}{[DATE]}
\cvcert{[CERTIFICATION]}{[DATE]}

% ---------- References ----------
\section{References}

oindent	extbf{[REFERENCE_NAME]} -- [TITLE], [COMPANY] -- \href{mailto:[REFERENCE_EMAIL]}{[REFERENCE_EMAIL]}\par

\end{document}
